%%% BEGIN fig:dis-oneloop
\begin{figure}[htbp]
\centering
% photons, photon vertices, cut and external labels common to the cut graphs (a)-(d)
\newcommand{\toneDIS}{%
  \draw[photon] (-1.5,2.6) -- (-1.0,1.8);
  \draw[photon] (1.0,1.8) -- (1.5,2.6);
  \node[vertex] at (-1.0,1.8) {}; \node[vertex] at (1.0,1.8) {};
  \draw[cut] (0,-0.25) -- (0,2.65);
  \node[below] at (-1.0,0) {$p$}; \node[below] at (1.0,0) {$p$};
  \node[left,inner sep=2pt] at (-1.5,2.6) {$q$};}
\begin{subfigure}[b]{0.235\textwidth}\centering
\begin{tikzpicture}[fd]
  \draw[quark] (-1.0,0) -- (-1.0,0.75);
  \draw[quark] (-1.0,0.75) -- (-1.0,1.8);
  \draw[quark] (-1.0,1.8) -- (-0.35,1.8);
  \draw[quark=0.65] (-0.35,1.8) -- (1.0,1.8);
  \draw[quark] (1.0,1.8) -- (1.0,0);
  \draw[gluon] (-1.0,0.75) to[out=0,in=-90] (-0.35,1.8);
  \toneDIS
  \node[vertex] at (-1.0,0.75) {}; \node[vertex] at (-0.35,1.8) {};
\end{tikzpicture}
\caption{虚修正（及其镜像）}
\end{subfigure}\hfill
\begin{subfigure}[b]{0.235\textwidth}\centering
\begin{tikzpicture}[fd]
  \draw[quark] (-1.0,0) -- (-1.0,0.75);
  \draw[quark] (-1.0,0.75) -- (-1.0,1.8);
  \draw[quark=0.72] (-1.0,1.8) -- (1.0,1.8);
  \draw[quark] (1.0,1.8) -- (1.0,0.75);
  \draw[quark] (1.0,0.75) -- (1.0,0);
  \draw[gluon] (-1.0,0.75) -- (1.0,0.75);
  \toneDIS
  \node[vertex] at (-1.0,0.75) {}; \node[vertex] at (1.0,0.75) {};
\end{tikzpicture}
\caption{初态辐射}
\end{subfigure}\hfill
\begin{subfigure}[b]{0.235\textwidth}\centering
\begin{tikzpicture}[fd]
  \draw[quark] (-1.0,0) -- (-1.0,1.8);
  \draw[quark] (-1.0,1.8) -- (-0.5,1.8);
  \draw[quark=0.25] (-0.5,1.8) -- (0.5,1.8);
  \draw[quark] (0.5,1.8) -- (1.0,1.8);
  \draw[quark] (1.0,1.8) -- (1.0,0);
  \draw[gluon] (-0.5,1.8) arc[start angle=180,end angle=0,radius=0.5];
  \toneDIS
  \node[vertex] at (-0.5,1.8) {}; \node[vertex] at (0.5,1.8) {};
\end{tikzpicture}
\caption{末态辐射}
\end{subfigure}\hfill
\begin{subfigure}[b]{0.235\textwidth}\centering
\begin{tikzpicture}[fd]
  \draw[quark] (-1.0,0) -- (-1.0,0.75);
  \draw[quark] (-1.0,0.75) -- (-1.0,1.8);
  \draw[quark=0.3] (-1.0,1.8) -- (0.5,1.8);
  \draw[quark] (0.5,1.8) -- (1.0,1.8);
  \draw[quark] (1.0,1.8) -- (1.0,0);
  \draw[gluon] (-1.0,0.75) to[out=0,in=-90] (0.5,1.8);
  \toneDIS
  \node[vertex] at (-1.0,0.75) {}; \node[vertex] at (0.5,1.8) {};
\end{tikzpicture}
\caption{干涉项（及其镜像）}
\end{subfigure}

\bigskip
\begin{subfigure}[b]{0.4\textwidth}\centering
\begin{tikzpicture}[fd]
  \coordinate (A) at (0,-0.65); \coordinate (B) at (0,0.65);
  \draw[gluon] (-1.4,-0.65) -- (A);
  \draw[photon] (-1.4,0.65) -- (B);
  \draw[quark] (1.4,-0.65) -- (A);
  \draw[quark] (A) -- (B);
  \draw[quark] (B) -- (1.4,0.65);
  \node[vertex] at (A) {}; \node[vertex] at (B) {};
  \node[left] at (-1.4,-0.65) {$g(p)$};
  \node[left] at (-1.4,0.65) {$\gamma^*(q)$};
  \node[right] at (1.4,0.65) {$q(k_1)$};
  \node[right] at (1.4,-0.65) {$\bar q(k_2)$};
\end{tikzpicture}
\caption{$\gamma^*g\to q\bar q$，$t$ 道}
\end{subfigure}\hspace{1em}
\begin{subfigure}[b]{0.4\textwidth}\centering
\begin{tikzpicture}[fd]
  \coordinate (A) at (0,-0.65); \coordinate (B) at (0,0.65);
  \draw[gluon] (-1.4,-0.65) -- (A);
  \draw[photon] (-1.4,0.65) -- (B);
  \draw[quark] (1.4,0.65) -- (B);
  \draw[quark] (B) -- (A);
  \draw[quark] (A) -- (1.4,-0.65);
  \node[vertex] at (A) {}; \node[vertex] at (B) {};
  \node[left] at (-1.4,-0.65) {$g(p)$};
  \node[left] at (-1.4,0.65) {$\gamma^*(q)$};
  \node[right] at (1.4,0.65) {$\bar q(k_2)$};
  \node[right] at (1.4,-0.65) {$q(k_1)$};
\end{tikzpicture}
\caption{$\gamma^*g\to q\bar q$，$u$ 道}
\end{subfigure}
\caption{DIS 在 $\order{\as}$ 的部分子图。(a)--(d) 为 $\gamma^*q$ 道的割图（左侧为振幅，右侧为共轭振幅）：
(a) 一圈顶点修正与树图的干涉（外线自能在维数正规化中无标度，未画出）；(b) 初态辐射的平方；(c) 末态辐射的平方；
(d) 初态辐射与末态辐射的干涉。(e)、(f) 为玻色子—胶子聚合 $\gamma^*g\to q\bar q$ 的两个振幅图（虚光子分别与夸克、反夸克耦合，内线传播子分别为 $1/\hat t$ 与 $1/\hat u$，$\hat t=(q-k_1)^2$，$\hat u=(q-k_2)^2$），
它们的平方与干涉给出胶子道系数函数。}
\label{fig:dis-oneloop}
\end{figure}
%%% END fig:dis-oneloop
